\documentclass{jsarticle}
\usepackage{ascmac,amsmath,amssymb,amsthm,bm,physics,arydshln,mathcomp}
\newtheorem*{question*}{問}
\newtheorem*{answer*}{解答}
\begin{document}
\begin{question*}
\end{question*}

    Let $P_n$, be the $(n + 1)$-dimensional vector space over $\mathbb{C}$ consisting of all polynomials in $x$ of degree at most $n$ with coefficients in $\mathbb{C}$. Let $f_n$ : $P_n \to P_n$ be the linear transformation given by 

\vspace{10pt}
\begin{center}
$f_n : p(x) \mapsto p(x) + p''(x) + p'''(x), \ p(x) \in P_n$
\end{center}
\vspace{10pt}

    where $''(x)$ and $p'''(x)$ are the second and third derivatives of $p(x)$, respectively. 

\vspace{10pt}

    (1) Find the representation matrix of $f_4$; with respect to the basis$x^4, x^3, x^2, x, 1$ of $P_4$.   

\vspace{10pt}


    (2) Find all eigenvalues of $f_4$ and the corresponding eigenspaces. 

\vspace{10pt}

    (3) Find the minimal polynomial of $f_4$, as well as a Jordan canonical form of a representation matrix

\vspace{10pt}

    \ \ \ \ of $f_4$. (Here, there is no need to find an invertible 
    matrix that transforms a representation matrix 

\vspace{10pt}

   \ \ \ \ into a Jordan canonical form. Assume the same for (4).) 

\vspace{10pt}

    (4) For each $n \ge 1$, find a Jordan canonical form of a representation matrix of $f_n$. 

\vspace{10pt}

\par

\begin{answer*}
\end{answer*}

(1) \begin{align*}
(f(x^4), \ f(x^3), \ f(x^2), \ f(x), \ f(1)) 
  &= (x^4 + 12x^2 + 24x, \ x^3 + 6x + 6, \ x^2 + 2, \ x, \ 1) \\
  &=(x^4, \ x^3, \ x^2, \ x, \ 1) \begin{pmatrix} 1&0&0&0&0 \\ 0&1&0&0&0 \\ 12&0&1&0&0 \\ 24&6&0&1&0 \\ 0&6&2&0&1 \end{pmatrix}
\end{align*}

(2) 表現行列を$A$ とすると, $| tE - A | = \begin{vmatrix} t-1&0&0&0&0 \\ 0&t-1&0&0&0 \\ -12&0&t-1&0&0 \\ -24&-6&0&t-1&0 \\ 0&-6&-2&0&t-1 \end{vmatrix} = (t-1)^5$ なので固有値$1$. 

\ \ \ \ 固有ベクトルは, $E - A = \begin{pmatrix} 0&0&0&0&0 \\ 0&0&0&0&0 \\ -12&0&0&0&0 \\ -24&-6&0&0&0 \\ 0&-6&-2&0&0 \end{pmatrix} \to \begin{pmatrix} 1&0&0&0&0 \\ 0&1&0&0&0 \\ 0&0&1&0&0 \\ 0&0&0&0&0 \\ 0&0&0&0&0 \end{pmatrix}$ なので 

\ \ \ \ $c_1 \begin{pmatrix} 0\\0\\0\\1\\0 \end{pmatrix} + c_2 \begin{pmatrix} 0\\0\\0\\0\\1 \end{pmatrix} \ (c_1, c_2 \in \mathbb{R} - \{0\})$ となる. 

(3) $(A-E)^2 = \begin{pmatrix} 0&0&0&0&0 \\ 0&0&0&0&0 \\ 0&0&0&0&0 \\ 0&0&0&0&0 \\ 24&0&0&0&0 \end{pmatrix}$, \ $(A-E)^3 = O$ なので最小多項式は$(t-1)^3$. よって$f_4$ のJordan

\ \ \ \ 標準形を構成するJordan細胞の最大次数は$3$であり, さらに固有空間$W(1;A)$ の次元は$2$ であること

\ \ \ \ から固有値$1$ のJordan細胞は$2$個存在する. よって$f_4$ のJordan標準形は$J(1;3) \oplus J(1;2)$.

(4) $f_n$ の基底$\{x^n, \ x^{n-1}, \dotsb\dotsb, \ x, \ 1\}$ を考えると

$(f_n(x^n), \ f_n(x^{n-1}), \ \dotsb\dotsb, f_n(x), \ f_n(1))$ 

= $(x^n, \ x^{n-1}, \cdots\cdots, x, \ 1)
\begin{pmatrix} 1 & 0 & 0 & 0 & \dotsb & 0 & 0 & 0 \\
                0 & 1 & 0 & 0 & \dotsb & 0 & 0 & 0 \\
                n(n-1) & 0 & 1 & 0 & \dotsb & 0 & 0 & 0 \\
                n(n-1)(n-2) & (n-1)(n-2) & 0 & 1 & \dotsb & 0 & 0 & 0 \\
                \vdots & \vdots & \vdots & \vdots & \ddots & \vdots & \vdots & \vdots \\
                0 & 0 & 0 & 0 & \cdots & 1 & 0 & 0 \\
                0 & 0 & 0 & 0 & \cdots & 0 & 1 & 0 \\
                0 & 0 & 0 & 0 & \cdots & 2 & 0 & 1 \\
                \end{pmatrix}$

となり, $f_n$ の表現行列を$A$ とすると

$(A-E)^{n-2} \neq O, \ (A-E)^{n-1} = O $ となることがわかる. 
(厳密な議論をするなら帰納法？)

よって最小多項式は$(t-1)^{n-1}$. そして, 

\begin{align*}
A-E &= \begin{pmatrix} 0 & 0 & 0 & 0 & \dotsb & 0 & 0 & 0 \\
                0 & 0 & 0 & 0 & \dotsb & 0 & 0 & 0 \\
                n(n-1) & 0 & 0 & 0 & \dotsb & 0 & 0 & 0 \\
                n(n-1)(n-2) & (n-1)(n-2) & 0 & 0 & \dotsb & 0 & 0 & 0 \\
                \vdots & \vdots & \vdots & \vdots & \ddots & \vdots & \vdots & \vdots \\
                0 & 0 & 0 & 0 & \cdots & 0 & 0 & 0 \\
                0 & 0 & 0 & 0 & \cdots & 0 & 0 & 0 \\
                0 & 0 & 0 & 0 & \cdots & 2 & 0 & 0 \\
                \end{pmatrix} \\
    &\to \begin{pmatrix} 1 & 0 & 0 & 0 & \dotsb & 0 & 0 & 0 \\
                0 & 1 & 0 & 0 & \dotsb & 0 & 0 & 0 \\
                0 & 0 & 1 & 0 & \dotsb & 0 & 0 & 0 \\
                0 & 0 & 0 & 1 & \dotsb & 0 & 0 & 0 \\
                \vdots & \vdots & \vdots & \vdots & \ddots & \vdots & \vdots & \vdots \\
                0 & 0 & 0 & 0 & \cdots & 2 & 0 & 0 \\
                0 & 0 & 0 & 0 & \cdots & 0 & 0 & 0 \\
                0 & 0 & 0 & 0 & \cdots & 0 & 0 & 0 \\
                \end{pmatrix} \\
\end{align*}

となるので固有空間$W(1;A)$ の次元は$2$.

したがって$f_n$ のJordan標準形は$J(1;n-1) \oplus J(1;2)$.




\end{document}
